\section{Finite sampling, residue curvature and the original control form}
\label{sec:periodized-curvature-control}

We retain the polynomial source, remainder map, coefficient basis and
arithmetic inner product from the periodized-source construction. The
following calculation propagates its proved finite-source errors to the
residue insertion and the original Laplacian. The quotient curvature
calculation also incorporates the exact Schur-complement argument supplied
by the coordinating Zeta task. Fix the original tensor degree
$k\in\mathbb Z_{\ge1}$ throughout the source construction.

Let $E=\mathbb C[S]/(\chi)$, with the original monic $\chi$ of degree
$q\ge1$, $A=M_S$, $e_0=[1]$, $\ell[P]=[S^{q-1}]\operatorname{rem}_\chi P$,
and $R=e_0\ell$. For $N\ge q-1$, let $\mathcal P_N$ be the original
polynomial coefficient space, $J_N:\mathcal P_N\twoheadrightarrow E$
the full remainder map, and $M_N>0$ its original source Gram. Keep
\[
 G=G_N=(J_NM_N^{-1}J_N^*)^{-1}.
 \tag{FC1}
\]
All adjoints denoted by $*$ are conjugate transposes in these specified
coefficient frames. For a positive form $T$, write
$B^{\dagger_T}=T^{-1}B^*T$ and $\|B\|_T$ for its induced operator norm.

\subsection{The actual finite sampling error}

Fix the common source degree $D\ge N$, $a>0$, integer $p\ge1$, circle
length $L>0$ and integer cutoff $J\ge1$. In the original diagonal log
coordinates $r=\log x_k$, $z_i=\log x_i-\log x_k$, retain
\[
 (\mathcal U_kF)(r,z)=e^{(kr+\sum_i z_i)/2}
 F(e^{r+z_1},\ldots,e^{r+z_{k-1}},e^r),\qquad
 \mathcal K=L^2(\mathbb R^{k-1},dz).
 \tag{FC2}
\]
The Jacobian is $d^kx=e^{kr+\sum_i z_i}\,dr\,dz$, proving directly that
this map preserves the original norm. Its derivative satisfies
$\mathcal U_kD^{(k)}=(-\partial_r+k/2)\mathcal U_k$.
Let $\Psi_D$ be the image of the degree-$D$ source columns. Define
\[
\begin{gathered}
 M_{a,\pm,D}=(e^{\pm ar}\Psi_D)^*(e^{\pm ar}\Psi_D),\qquad
 H_{p,D}=\int_{\mathbb R}(1+r^2)(\partial_r^p\Psi_D)^*
                  (\partial_r^p\Psi_D)\,dr,\\
 \kappa_{a,D}=\max_{c\ne0}
 \frac{c^*(M_{a,+,D}+M_{a,-,D})c}{c^*M_Dc},\qquad
 \lambda_{p,D}=\max_{c\ne0}\frac{c^*H_{p,D}c}{c^*M_Dc},\\
 \eta_P=\frac{\kappa_{a,D}}{e^{aL}-1},\qquad
 \eta_T=\frac{\lambda_{p,D}}{2p-1}
                   \left(\frac{L}{2\pi J}\right)^{2p-1},\qquad
 \delta=\eta_P+\eta_T .
\end{gathered}\tag{FC3}
\]
These are finite observations of the original source; the weighted forms
retain $x_k^{\pm2a}$ and $1+(\log x_k)^2$ in the original coordinates.
For positive budgets $\varepsilon_P,\varepsilon_T$ with
$\varepsilon_P+\varepsilon_T<1$, choose the literal values of (P64):
\[
\begin{gathered}
 L\ge a^{-1}\log(1+\kappa_{a,D}/\varepsilon_P),\\
 J\ge\max\left\{1,D,\frac{L}{2\pi}
       \left(\frac{\lambda_{p,D}}{(2p-1)\varepsilon_T}
       \right)^{1/(2p-1)}\right\},\qquad J\in\mathbb Z.
\end{gathered}\tag{FC3a}
\]
Substitution into (FC3) proves $\eta_P\le\varepsilon_P$ and
$\eta_T\le\varepsilon_T$, including $\lambda_{p,D}=0$.
Thus $0\le\delta<1$ has been obtained from actual finite quantities.
No estimate uniform in $h,k,D$ is inserted.

Use the plus Fourier transform $\widehat\psi(u)=\int\psi(r)e^{iur}dr$.
Let $M_N(L,J)$ be the full source Gram
\[
 M_N(L,J)_{ab}=\frac1L\sum_{|n|\le J}
 \left\langle\widehat{\Psi_N e_a}(2\pi n/L),
                  \widehat{\Psi_N e_b}(2\pi n/L)\right\rangle_{\mathcal K}.
 \tag{FC4}
\]
The $n=0$ term is included with its original factor $L^{-1}$.
For clarity, the estimates producing (FC3) can be checked without an
inverse operation. Unfolding periodization gives
$M_N(L)-M_N=\sum_{m\ne0}\mathcal C_N(mL)$. Weighted Cauchy--Schwarz gives
$|c^*\mathcal C_N(s)c|\le e^{-a|s|}
\sqrt{c^*M_{a,-,N}c\,c^*M_{a,+,N}c}$.
Summing both lattice directions bounds its absolute value by
$c^*(M_{a,-,N}+M_{a,+,N})c/(e^{aL}-1)$.
Also integration by parts and $\int(1+r^2)^{-1}dr=\pi$ give
$\|\widehat{\Psi_Nc}(u)\|^2\le\pi|u|^{-2p}c^*H_{p,N}c$.
Summing the two tails $|n|>J$, with
$\sum_{n>J}n^{-2p}\le J^{1-2p}/(2p-1)$, yields exactly
the coefficient of $\lambda_{p,D}$ in (FC3). Restriction from the
degree-$D$ source preserves all these quadratic-form bounds. Therefore
\[
 (1-\eta_P-\eta_T)M_N\preceq M_N(L,J)
       \preceq(1+\eta_P)M_N.
 \tag{FC5}
\]
The asymmetry records that the omitted Fourier tail is positive.
Put $\alpha=1-\eta_P-\eta_T>0$, $\beta=1+\eta_P$, and retain
\[
 H=G_N(L,J)=(J_NM_N(L,J)^{-1}J_N^*)^{-1}.
 \tag{FC6}
\]
Completing the square proves that $v^*Gv$ is the minimum of $c^*M_Nc$
over $J_Nc=v$, and the same statement holds for $H$. Taking these
minima in (FC5) proves
\[
 \alpha G\preceq H\preceq\beta G,
 \qquad (1-\delta)G\preceq H\preceq(1+\delta)G.
 \tag{FC7}
\]
The coefficient minimizers are $C_N=M_N^{-1}J_N^*G$ and
$C_N(L,J)=M_N(L,J)^{-1}J_N^*H$. Their difference belongs to
$\chi\mathcal P_{N-q}$. Fixed monic division writes its polynomial
as $\chi Q=\sum_i h(s_i)Q_i(\mathbf s)$. The original tensor primitive
replaces $F_h$ by $\phi_*$ in slot $i$, applies $Q_i(D_1,\ldots,D_k)$,
and has coefficient $(-1)^{i-1}$. The tensor differential contributes
the same sign. Thus this change of representative has an explicit
primitive with every original summand retained.

\subsection{A dimension-free relative estimate for the residue coefficient}

Let $I:F\hookrightarrow E$ be an actual nonzero proper $A$-invariant
constituent and set $L_F=\ell I$. For every positive form $T$ define
\[
\begin{gathered}
 T_F=I^*TI,\qquad P_{T,F}=I T_F^{-1}I^*T,\\
 a_F(T)=\min_{x\in F}(e_0-Ix)^*T(e_0-Ix)
       =e_0^*T(1-P_{T,F})e_0,\\
 b_F(T)=L_F T_F^{-1}L_F^*,\qquad
 \mathfrak c_F(T)=a_F(T)b_F(T),\qquad
 \sigma(T)^2=(e_0^*Te_0)(\ell T^{-1}\ell^*).
\end{gathered}\tag{FC8}
\]
Here $q>1$. To verify strict positivity, $I(F)$ is an ideal of $E$
because it is invariant under every polynomial in $A$. Its polynomial
preimage is $(g)$ with $\chi=gh$, and the original ideal basis is
$[g],\ldots,[S^{p_F-1}g]$, where $p_F=\deg h=\dim F$.
Its last member is monic of degree $q-1$, so $L_F\ne0$; properness gives
$e_0\notin I(F)$. These statements prove $a_F(T),b_F(T)>0$ without
discarding repeated factors.

The same constrained-minimum argument as in (FC7) gives
$\alpha a_F(G)\le a_F(H)\le\beta a_F(G)$.
Compression and inverse order give
$\beta^{-1}b_F(G)\le b_F(H)\le\alpha^{-1}b_F(G)$.
Inverse order follows, for example, from the exact identity
$w^*T^{-1}w=\max_x(2\operatorname{Re}w^*x-x^*Tx)$.
Multiplication proves the full coefficient estimate
\[
 \boxed{\frac{\alpha}{\beta}\le
       \frac{\mathfrak c_F(H)}{\mathfrak c_F(G)}
           \le\frac{\beta}{\alpha},\qquad
 \left|\log\frac{\mathfrak c_F(H)}{\mathfrak c_F(G)}\right|
           \le\log\frac{\beta}{\alpha}.}
 \tag{FC9}
\]
Applying the same two form bounds to $e_0$ and the inverse form to
$\ell^*$ also proves
\[
 \frac{\alpha}{\beta}\sigma(G)^2\le\sigma(H)^2
       \le\frac{\beta}{\alpha}\sigma(G)^2.
 \tag{FC10}
\]
Every number is evaluated in its stated original or sampled metric.
Neither (FC9) nor (FC10) has a dimension factor.

For the actual period matrix $\Pi(u,t)$, $u\ne0$, put
$T_{\rm per}(t)=\Pi(t)^*\Pi(t)$ and let $m_H(t),M_H(t)$ be the least
and greatest eigenvalues of $H^{-1}T_{\rm per}(t)$. Their positivity
follows from the nonzero full period determinant. Applying the same
calculation, now between $H$ and $T_{\rm per}$, gives
\[
 \frac{m_H(t)}{M_H(t)}\mathfrak c_F(H)
 \le\mathfrak c_F(T_{\rm per}(t))
 \le\frac{M_H(t)}{m_H(t)}\mathfrak c_F(H).
 \tag{FC11}
\]
The period equation $\Pi'=-\Pi(A+tR)/u$ implies, with $Y=\Pi I$ and
$P=Y(Y^*Y)^{-1}Y^*$,
$(1-P)Y'=-t(1-P)\Pi e_0L_F/u$. Differentiating
$\log\det(Y^*Y)$ gives
$\partial_{\bar t}\partial_t\log\det(Y^*Y)
=\operatorname{Tr}((Y^*Y)^{-1}Y'^*(1-P)Y')$.
Consequently the precise connection to (FC9) is
\[
 \partial_{\bar t}\partial_t\log\det(Y^*Y)
       =\frac{|t|^2}{|u|^2}\mathfrak c_F(T_{\rm per}(t)).
 \tag{FC12}
\]
Equations (FC9)--(FC12) retain the actual comparison spectrum; they do
not substitute the sampled coefficient for the period coefficient.

There is also a useful exact endpoint bound. For two admitted degrees
$i\le j$, the original representative spaces increase, so $G_i\succeq
G_j>0$. Let $\theta$ and $\zeta_+$ be the least and greatest
eigenvalues of the actual relative operator $G_i^{-1}G_j$.
Then $\theta G_i\preceq G_j\preceq\zeta_+G_i$ with
$0<\theta\le\zeta_+\le1$. Applying the proved minimum and inverse
comparisons with these two factors gives the first inequality below.
All eigenvalues lie in $(0,1]$, so their product is at most their
least member. This proves the remaining inequalities and retains
the sharper endpoint coefficient from (RCX21):
\[
 \left|\log\frac{\mathfrak c_F(G_j)}{\mathfrak c_F(G_i)}\right|
 \le\log\frac{\zeta_+}{\theta}
 \le-\log\theta\le\log\frac{\det G_i}{\det G_j}.
 \tag{FC13}
\]
The last inequality retains the correct direction. It does not infer
an upper volume estimate from a scalar curvature observation.
For four original endpoints, (FC7) similarly gives the fully finite
bound $|\mathcal B(L,J)-\mathcal B|\le2q\log(\beta/\alpha)$.

\subsection{Transporting the control form through the sampling error}

The control also depends on the unchanged operator, so we retain its
complete error term. For $t\in\mathbb C$, retaining that tensor degree
$k$, set
\[
\begin{gathered}
 A_t=A+tR,\quad T_t=A_t-\tfrac k2I,\quad
 W_G(t)=A_t^*G+GA_t-kG,\quad X_G(t)=G^{-1}W_G(t),\\
 W_H(t)=A_t^*H+HA_t-kH,\quad X_H(t)=H^{-1}W_H(t),\quad
 B=G^{-1}(H-G).
\end{gathered}\tag{FC14}
\]
The matrix identities in this subsection and in (FC24)--(FC29)
also hold for any real scalar $k$ when considered on the displayed
finite coefficient space; the source realization retains its given
integer tensor degree.
Then $B^{\dagger_G}=B$, $\|B\|_G\le\delta$, and $H=G(I+B)$.
No coordinate vector or inner product has been rescaled. Direct
multiplication proves the exact identity
\[
 (I+B)X_H(t)=X_G(t)+T_t^{\dagger_G}B+BT_t.
 \tag{FC15}
\]
Let $\epsilon_G(t)=\|X_G(t)\|_G$, $\epsilon_H(t)=\|X_H(t)\|_H$,
and $\tau_G(t)=\|T_t\|_G$. The two $X$ operators are self-adjoint
in their respective metrics, so their norms are the suprema of their
absolute quadratic forms divided by the respective squared norms.
For every original $v$,
\[
 |v^*(W_H(t)-W_G(t))v|
 =|\langle T_tv,Bv\rangle_G+\langle Bv,T_tv\rangle_G|
 \le2\delta\tau_G(t)\,v^*Gv.
 \tag{FC16}
\]
Using (FC7) in both directions proves
\[
 \boxed{\epsilon_H(t)\le
  \frac{\epsilon_G(t)+2\delta\tau_G(t)}{\alpha},\qquad
 \epsilon_G(t)\le\beta\epsilon_H(t)+2\delta\tau_G(t).}
 \tag{FC17}
\]
Thus the control transfer retains the actual generator size as well as
the relative metric error. Both quantities are finite observations of
the original matrices. The exact LC additional radii
$|t|\sigma(G)$ for $q>1$, and $2|\operatorname{Re}t|$ for $q=1$,
remain available through (FC10).

For completeness, define the evaluated upper bound
$e=\beta\epsilon_H(t)+2\delta\tau_G(t)$ and $L_t=A_t^2-kA_t$.
Expansion gives $GL_t=W_G(t)T_t-T_t^*GT_t-k^2G/4$.
For each nonzero original vector $v$, put
$z=v^*GL_tv/(v^*Gv)$, $r=\|T_tv\|_G/\|v\|_G$ and
$a+ib=\langle v,X_G(t)T_tv\rangle_G/(v^*Gv)$.
Equation (FC17) gives $a^2+b^2\le e^2r^2$, while the expansion gives
$\operatorname{Re}z=a-r^2-k^2/4$ and $\operatorname{Im}z=b$.
The identities $er-r^2=e^2/4-(r-e/2)^2$ and
\[
 e^2\left(\frac{e^2-k^2}{4}-\operatorname{Re}z\right)-b^2
 =(a-e^2/2)^2+e^2r^2-a^2-b^2
\]
prove the actual original-metric enclosure
\[
 \operatorname{Re}z\le\frac{e^2-k^2}{4},\qquad
 (\operatorname{Im}z)^2\le e^2
       \left(\frac{e^2-k^2}{4}-\operatorname{Re}z\right).
 \tag{FC18}
\]
No division by $e$ occurs, so the zero-radius case is included.

\subsection{The finite Fourier boundary and the quotient compensation}

Let $Q_J$ be the orthogonal Fourier projection to $|n|\le J$, with
the full $\mathcal K$ fibre at each retained mode. It commutes with
$D_L=-\partial_r+k/2$. Use the actual least-norm representative for
$M_N(L,J)$ and write
\[
\begin{gathered}
 r^{L,J}=\mathcal V_{h,k}C_N(L,J),\qquad
 \mathscr R=Q_J\mathcal P_L\mathcal U_kr^{L,J},\\
 \mathscr B=Q_J\mathcal P_L\mathcal U_k
                (D^{(k)}r^{L,J}-r^{L,J}A).
\end{gathered}\tag{FC19}
\]
The explicit polynomial-division primitive following (FC7) applies
to this relation as well. Direct intertwining gives
$D_L\mathscr R-\mathscr RA=\mathscr B$ and $\mathscr R^*\mathscr R=H$.
Since $D_L^*+D_L=k$ on these smooth finite-mode columns, multiplication
and then one further application of $D_L$ prove
\[
\begin{gathered}
 W_H(0)=-(\mathscr R^*\mathscr B+\mathscr B^*\mathscr R),\\
 (D_L^2-kD_L)\mathscr R-\mathscr R(A^2-kA)
                    =(D_L-k)\mathscr B+\mathscr BA.
\end{gathered}\tag{FC20}
\]
In particular (FC17) transports an actual boundary control computed
from (FC20) to the original Gram. Finite mode count does not assign
interval enclosures to its arithmetic coefficients; such enclosures
have to be proved for any numerical application.

Finally retain the exact coefficient quotient $E/I(F)$ and a constant
lift matrix $J_F$ of its fixed basis. Then $C=[I,J_F]$ is invertible.
For $T_{\rm per}=\Pi^*\Pi$ the least-period-norm lift and its Gram are
\[
\begin{gathered}
 Q(t)=J_F-I(I^*T_{\rm per}I)^{-1}I^*T_{\rm per}J_F,\\
 K(t)=J_F^*T_{\rm per}J_F
   -J_F^*T_{\rm per}I(I^*T_{\rm per}I)^{-1}I^*T_{\rm per}J_F.
\end{gathered}\tag{FC21}
\]
Completing the square in $\|\Pi(J_Fv+Ix)\|^2$ proves uniqueness,
positivity, $\pi Q=\pi J_F$ and $I^*T_{\rm per}Q=0$.
The change $J_F\mapsto J_F+IB_0$ cancels exactly in $Q$ and $K$;
a quotient basis change $J_F\mapsto J_FD_0$ gives $K\mapsto D_0^*KD_0$.
Block elimination in $C^*T_{\rm per}C$ gives
\[
 \det(I^*T_{\rm per}I)\det K
       =|\det C|^2|\det\Pi|^2.
 \tag{FC22}
\]
The complete period determinant is a nonzero constant times a
holomorphic exponential in $t$ at fixed $u$. Its retained squared
constant is $(2\pi|u|)^q$; hence its log modulus squared has zero mixed
curvature. Differentiating (FC22) and using (FC12) proves
\[
 \boxed{\partial_{\bar t}\partial_t\log\det K(t)
       =-\frac{|t|^2}{|u|^2}\mathfrak c_F(T_{\rm per}(t)).}
 \tag{FC23}
\]
Both actual curvatures vanish at $t=0$, and their nonzero leading
coefficients have opposite signs. More precisely applying
$\partial_t\partial_{\bar t}$ to (FC12) at zero gives
$\partial_t^2\partial_{\bar t}^2\log\det(I^*T_{\rm per}I)|_0
=\mathfrak c_F(T_{\rm per}(0))/|u|^2$; the same derivative of
$\log\det K$ is its negative. The respective $t^2\bar t^2$ Taylor
coefficients are these numbers divided by $2!2!$.

Apply the same Schur formula to the unchanged $G$ and sampled $H$ to
obtain $K_G,K_H$. The fibre minima in the same quotient give
$\alpha K_G\preceq K_H\preceq\beta K_G$.
The literal comparison endomorphism $K_G^{-1}K(t)$ satisfies
$K_G(K_G^{-1}K(t))=K(t)$ and has log determinant curvature (FC23).
This supplies the metric morphism on the original quotient, together
with its finite sampling error and its exact curvature compensation.
No coefficient class, generalized eigenspace or support label is
removed by these identities.

\subsection{The exact commutator allowance for the same two metrics}

The coordinating Zeta task supplied the following sharper transfer,
which we prove in full for the actual $G,H,A_t$ already defined.
It retains the specific metric discrepancy that interacts with the
generator. All adjoints in this subsection are for $G$.
Put $\mathcal C=G^{-1}H=I+B$, using the same $B$ as in (FC14).
This is a positive $G$-self-adjoint endomorphism with spectrum in
$[\alpha,\beta]$. This symbol denotes the metric endomorphism;
the constant quotient frame $C=[I,J_F]$ in (FC22) remains unchanged.
Subtract $\mathcal C X_G(t)$ from (FC15), in the displayed order:
\[
 \boxed{X_H(t)-X_G(t)
       =\mathcal C^{-1}[T_t^{\dagger_G},B].}
 \tag{FC24}
\]
Indeed $X_G=T_t^{\dagger_G}+T_t$, so the numerator of this
difference is $T_t^{\dagger_G}B-BT_t^{\dagger_G}$.

Let $\mathcal S=\mathcal C^{1/2}$ be the unique positive
$G$-self-adjoint square root. The finite spectral theorem constructs
it by replacing each positive eigenvalue of the actual $\mathcal C$
by its positive square root. Uniqueness follows because a positive
square root commutes with its square and, on each eigenspace of that
square, positivity forces its eigenvalues to be that positive root.
Thus no additional metric is selected. The exact map
$\mathcal S:(E,H)\to(E,G)$ is an onto isometry, since
$\mathcal S^*G\mathcal S=G\mathcal S^2=G\mathcal C=H$.
Its inverse is the actual $\mathcal S^{-1}$.
In this comparison set
\[
\begin{gathered}
 Y_H=\mathcal S X_H\mathcal S^{-1}
       =\mathcal S^{-1}T_t^{\dagger_G}\mathcal S
                       +\mathcal S T_t\mathcal S^{-1},\\
 Z=[\mathcal S,T_t]\mathcal S^{-1},\qquad
 Y_H-X_G=Z+Z^{\dagger_G}.
\end{gathered}\tag{FC25}
\]
The adjoint of $Z$ is
$\mathcal S^{-1}[T_t^{\dagger_G},\mathcal S]$; inserting it
proves the last identity with both ordered terms present.
The isometry gives $\|Y_H\|_G=\epsilon_H(t)$, so the reverse
triangle inequality gives
\[
 |\epsilon_H(t)-\epsilon_G(t)|
 \le\|Z+Z^{\dagger_G}\|_G
 \le2\|[\mathcal S,T_t]\mathcal S^{-1}\|_G.
 \tag{FC26}
\]

The commutator $Y=[\mathcal S,T_t]$ satisfies the Sylvester equation
$\mathcal S Y+Y\mathcal S=[\mathcal C,T_t]=[B,T_t]$.
Its exact solution is
\[
 [\mathcal S,T_t]=\int_0^\infty
       e^{-r\mathcal S}[B,T_t]e^{-r\mathcal S}\,dr.
 \tag{FC27}
\]
To prove convergence, the original $G$-operator norm satisfies
$\|e^{-r\mathcal S}\|_G\le e^{-r\sqrt\alpha}$.
Differentiating the integrand and integrating its two terms gives
$\mathcal S Y+Y\mathcal S=[B,T_t]$, with value $[B,T_t]$
at zero and limit zero at infinity. For uniqueness, in the spectral
decomposition of $\mathcal S$, each matrix block of a homogeneous
solution is multiplied by the positive sum of its two eigenvalues;
all such blocks must vanish. Therefore the integral equals the
stated commutator, rather than a separate solution chosen for a bound.
Its norm is at most $\|[B,T_t]\|_G/(2\sqrt\alpha)$, while
$\|\mathcal S^{-1}\|_G\le\alpha^{-1/2}$. Substitution proves
\[
\begin{gathered}
 \boxed{|\epsilon_H(t)-\epsilon_G(t)|
       \le\|Z+Z^{\dagger_G}\|_G
       \le\alpha^{-1}\|[B,T_t]\|_G,}\\
 [B,T_t]=[B,A]+t[B,R],\qquad
 [B,R]v=(Be_0)\ell(v)-e_0\ell(Bv).
\end{gathered}\tag{FC28}
\]
The scalar $k/2$ commutes exactly, and both residue terms remain.
In particular the fully evaluated number
\[
 e_{\rm comm}(t)=\epsilon_H(t)+\|Z+Z^{\dagger_G}\|_G
 \le\epsilon_H(t)+\alpha^{-1}\|[B,A]+t[B,R]\|_G
 \tag{FC29}
\]
is another valid original-metric control radius in (FC18).
Taking the lesser of this number and the radius there remains valid
because both have just been proved to bound $\epsilon_G(t)$.
If $[B,T_t]=0$, its $G$-adjoint gives
$[T_t^{\dagger_G},B]=0$, and uniqueness in (FC27) gives
$[\mathcal S,T_t]=0$. Equations (FC24)--(FC26) then prove
$X_H=X_G$, $Y_H=X_G$ and $\epsilon_H=\epsilon_G$ exactly.

These identities are covariant under every invertible coordinate
map $U:E'\to E$. Indeed
$G'=U^*GU$, $H'=U^*HU$, $T_t'=U^{-1}T_tU$ and
$B'=U^{-1}BU$ imply $\mathcal S'=U^{-1}\mathcal S U$ by
positive-square-root uniqueness in $G'$. Every commutator and
integrand above transforms by this same similarity, and every norm
is preserved by the stated coordinate map. Thus (FC28) is a
comparison of the two original source observations with their full
generator; its arithmetic commutator still requires evaluation for
any numerical or uniform endpoint application.
